\documentclass[aspectratio=169,10pt]{beamer}

% ==============================================================================
% Plantilla de Presentaciones - Universidad Icesi
% Guía de DevOps: Scripts, Docker, GitFlows y GitHub Actions
% ==============================================================================

\usepackage[utf8]{inputenc}
\usepackage[spanish]{babel}
\usepackage{tikz}
\usepackage{tcolorbox}
\usepackage{listings}
\usepackage{booktabs}
\usepackage{hyperref}
\usepackage{amsmath}

% Definición de Paleta de Colores Institucional ICESI
\definecolor{IcesiNavy}{HTML}{002855}
\definecolor{IcesiBlue}{HTML}{003366}
\definecolor{IcesiCyan}{HTML}{00A9E0}
\definecolor{IcesiOrange}{HTML}{F39200}
\definecolor{IcesiDark}{HTML}{0B132B}
\definecolor{IcesiLight}{HTML}{F4F7FB}
\definecolor{CodeBg}{HTML}{1E293B}

% Configuración del Tema Beamer
\usetheme{Madrid}
\usecolortheme{default}

% Personalización de Colores Beamer
\setbeamercolor{structure}{fg=IcesiBlue}
\setbeamercolor{title}{fg=white,bg=IcesiNavy}
\setbeamercolor{frametitle}{fg=white,bg=IcesiNavy}
\setbeamercolor{section in head/foot}{fg=white,bg=IcesiBlue}
\setbeamercolor{author in head/foot}{fg=white,bg=IcesiNavy}
\setbeamercolor{date in head/foot}{fg=white,bg=IcesiBlue}
\setbeamercolor{block title}{fg=white,bg=IcesiBlue}
\setbeamercolor{block body}{fg=IcesiDark,bg=IcesiLight}
\setbeamercolor{block title alerted}{fg=white,bg=IcesiOrange}
\setbeamercolor{block body alerted}{fg=IcesiDark,bg=IcesiLight}
\setbeamercolor{item}{fg=IcesiCyan}

% Configuración de Listings para Código Fuente
\lstset{
  backgroundcolor=\color{CodeBg},
  basicstyle=\ttfamily\scriptsize\color{white},
  keywordstyle=\color{IcesiCyan}\bfseries,
  commentstyle=\color{gray}\itshape,
  stringstyle=\color{IcesiOrange},
  numbers=none,
  breaklines=true,
  showstringspaces=false,
  tabsize=2,
  frame=single,
  rulecolor=\color{IcesiBlue}
}

% Datos de la Presentación
\title[DevOps \& CI/CD]{DevOps \& CI/CD para Ingeniería de Software}
\subtitle{Scripts, Docker Multi-Stage, GitFlows y GitHub Actions}
\author[IasLab · Icesi]{Departamento de TIC · Universidad Icesi}
\institute[ICESI]{Facultad de Ingeniería · Infraestructura IasLab}
\date{\today}

\begin{document}

% ------------------------------------------------------------------------------
% Diapositiva 1: Portada
% ------------------------------------------------------------------------------
\begin{frame}[plain]
  \titlepage
\end{frame}

% ------------------------------------------------------------------------------
% Diapositiva 2: Agenda
% ------------------------------------------------------------------------------
\begin{frame}{Ruta de Aprendizaje}
  \tableofcontents
\end{frame}

% ==============================================================================
\section{Módulo 1: Automatización con Scripts}
% ==============================================================================

\begin{frame}[fragile]{Módulo 1: Shell Scripting Defensivo}
  \begin{columns}
    \begin{column}{0.5\textwidth}
      \begin{block}{Directivas de Ejecución Segura}
        \begin{itemize}
          \item \textbf{\texttt{set -e}}: Detiene la ejecución ante cualquier error.
          \item \textbf{\texttt{set -u}}: Trata variables nulas como errores.
          \item \textbf{\texttt{set -o pipefail}}: Captura fallos en pipelines (\texttt{|}).
        \end{itemize}
      \end{block}
    \end{column}
    \begin{column}{0.5\textwidth}
      \begin{lstlisting}[language=bash,title=deploy-safe.sh]
#!/usr/bin/env bash
set -euo pipefail

TARGET="${1:-/home/iaslab/app}"

if [ ! -d "$TARGET" ]; then
  echo "Error: Directorio no existe"
  exit 1
fi

docker compose pull && docker compose up -d
      \end{lstlisting}
    \end{column}
  \end{columns}
\end{frame}

% ==============================================================================
\section{Módulo 2: Contenerización con Docker}
% ==============================================================================

\begin{frame}[fragile]{Módulo 2: Multi-Stage Builds en Docker}
  \begin{columns}
    \begin{column}{0.5\textwidth}
      \begin{block}{Backend: Spring Boot (Java 17)}
        \begin{lstlisting}[language=bash]
# 1. Build Stage
FROM eclipse-temurin:17-jdk AS builder
WORKDIR /app
COPY . .
RUN ./gradlew bootJar --no-daemon

# 2. Runtime Stage (< 150MB)
FROM eclipse-temurin:17-jre
USER spring
COPY --from=builder /app/build/*.jar app.jar
EXPOSE 9091
ENTRYPOINT ["java","-jar","/app.jar"]
        \end{lstlisting}
      \end{block}
    \end{column}
    \begin{column}{0.5\textwidth}
      \begin{block}{Frontend: React + Nginx Alpine}
        \begin{lstlisting}[language=bash]
# 1. Build Stage
FROM node:20-alpine AS builder
WORKDIR /app
COPY . .
RUN npm ci && npm run build

# 2. Runtime Stage (< 25MB)
FROM nginx:alpine
COPY --from=builder /app/dist /usr/share/nginx/html
EXPOSE 80
CMD ["nginx","-g","daemon off;"]
        \end{lstlisting}
      \end{block}
    \end{column}
  \end{columns}
\end{frame}

\begin{frame}[fragile]{Módulo 2: Orquestación y Reverse Proxy Nginx}
  \begin{columns}
    \begin{column}{0.5\textwidth}
      \begin{block}{docker-compose.yml}
        \begin{lstlisting}[language=bash]
services:
  web:
    image: github-action-test:latest
    container_name: test-app
    restart: unless-stopped
    ports:
      - "127.0.0.1:9088:80"
        \end{lstlisting}
      \end{block}
    \end{column}
    \begin{column}{0.5\textwidth}
      \begin{block}{Nginx Host Proxy (/home/iaslab/*/nginx.conf)}
        \begin{lstlisting}[language=bash]
location = /iaslab/app {
    return 301 /iaslab/app/$is_args$args;
}

location /iaslab/app/ {
    proxy_pass http://127.0.0.1:9088/;
    proxy_set_header Host $host;
    proxy_set_header X-Real-IP $remote_addr;
}
        \end{lstlisting}
      \end{block}
    \end{column}
  \end{columns}
\end{frame}

% ==============================================================================
\section{Módulo 3: Flujos de Trabajo en Git (GitFlow)}
% ==============================================================================

\begin{frame}{Módulo 3: Estrategia de Ramas y SemVer}
  \begin{center}
    \begin{tabular}{llll}
      \toprule
      \textbf{Rama} & \textbf{Propósito} & \textbf{Versionado SemVer} & \textbf{Ambiente} \\
      \midrule
      \texttt{main} & Producción Estable & \texttt{1.0.0} (Release) & Servidor Producción \\
      \texttt{release/*} & Pruebas QA / Staging & \texttt{1.0.0-RC1} & Servidor \texttt{grid100} \\
      \texttt{develop} & Integración Continua & \texttt{1.0.0-SNAPSHOT} & Local / Dev \\
      \texttt{feature/*} & Nuevas Funcionalidades & N/A & Ramas de Trabajo \\
      \texttt{hotfix/*} & Parches Críticos & \texttt{1.0.1} (Patch) & Producción Directa \\
      \bottomrule
    \end{tabular}
  \end{center}
  \vspace{0.3cm}
  \begin{alertblock}{Conventional Commits}
    \texttt{feat(auth): add JWT login} \quad $\vert$ \quad \texttt{fix(survey): resolve 500 null pointer} \quad $\vert$ \quad \texttt{ci: update runner}
  \end{alertblock}
\end{frame}

% ==============================================================================
\section{Módulo 4: CI/CD con GitHub Actions}
% ==============================================================================

\begin{frame}[fragile]{Módulo 4: Workflows Condicionales por Rama}
  \begin{columns}
    \begin{column}{0.5\textwidth}
      \begin{block}{Rama \texttt{main} (Despliegue Estático)}
        \begin{lstlisting}[language=bash]
on:
  push:
    branches: [ main ]

jobs:
  deploy-static:
    runs-on: [ self-hosted ]
    steps:
      - uses: actions/checkout@v4
      - run: |
          cp index.html /home/iaslab/app/
          sudo nginx -t && sudo systemctl reload nginx
        \end{lstlisting}
      \end{block}
    \end{column}
    \begin{column}{0.5\textwidth}
      \begin{block}{Rama \texttt{docker} (Contenedores)}
        \begin{lstlisting}[language=bash]
on:
  push:
    branches: [ docker ]

jobs:
  deploy-docker:
    runs-on: [ self-hosted ]
    steps:
      - uses: actions/checkout@v4
      - run: |
          docker compose build
          docker compose up -d
          sudo systemctl reload nginx
        \end{lstlisting}
      \end{block}
    \end{column}
  \end{columns}
\end{frame}

\begin{frame}{Módulo 4: Arquitectura Self-Hosted Runner (\texttt{grid100})}
  \begin{block}{¿Por qué usar Self-Hosted Runners en IasLab?}
    \begin{itemize}
      \item \textbf{Seguridad Saliente:} El runner se conecta a GitHub mediante \textit{HTTPS long-polling}; no requiere puertos públicos entrantes abiertos hacia el servidor.
      \item \textbf{Acceso a Recursos Internos:} Ejecuta comandos con acceso nativo al socket de Docker y a las bases de datos locales (PostgreSQL 18).
      \item \textbf{Sin Límites de Minutos:} Ejecuciones ilimitadas en infraestructura institucional.
    \end{itemize}
  \end{block}

  \begin{alertblock}{Monitoreo y URLs Activas}
    \begin{itemize}
      \item Portal del Runner: \url{https://pi2tools.icesi.edu.co/iaslab/github-action/}
      \item Despliegue Estático: \url{https://pi2tools.icesi.edu.co/iaslab/github-action-test/}
      \item Despliegue Docker: \url{https://pi2tools.icesi.edu.co/iaslab/github-action-docker/}
    \end{itemize}
  \end{alertblock}
\end{frame}

% ------------------------------------------------------------------------------
% Diapositiva Final
% ------------------------------------------------------------------------------
\begin{frame}[plain]
  \begin{center}
    \vspace{1cm}
    {\Huge \textbf{\color{IcesiBlue} Universidad Icesi}}\\[0.4cm]
    {\Large \textbf{\color{IcesiOrange} Departamento de TIC · Ingeniería de Software}}\\[1cm]
    {\large ¡Preguntas y Espacio de Laboratorio!}\\[0.5cm]
    \texttt{https://pi2tools.icesi.edu.co/iaslab/docs/}
  \end{center}
\end{frame}

\end{document}
